%! Author = lili
%! Date = 8/20/26


\newglossaryentry{wildtyp}{
    name={Wildtyp},
    description={In der Natur auftretende genetische Normalform
    des Genotyps und Phänotyps},
    type=defi
}

\newglossaryentry{transposon}{
    name={Transposon},
    description={Transposable Element - eine mobile Nukleinsäuresequenz (Mobiles genetisches Element) mit der Fähigkeit
    zur Transposition; charakterisiert durch umgebende repeats},
    type=defi
}


\newglossaryentry{transposition}{
    name={Transposition},
    description={Einbau eines
    DNA-Segments an einen anderen Ort im Genom.},
    type=defi
}
\newglossaryentry{tmd}{
    name={Transmembran-Domäne},
    description={Hydrophober Proteinabschnitt, der eine Zellmembran durchspannt und meist als Alpha-Helix vorliegt.
    Membranproteine können eine oder mehrere solcher Domänen besitzen.},
    type=defi
}

\newglossaryentry{tic}{
    name={translocon on the inner chloroplast membrane},
    description={Proteinkomplex in der inneren Chloroplastenmembran, der gemeinsam mit TOC den Import kernkodierter Proteine in den Chloroplasten vermittelt.},
    type=defi
}

\newglossaryentry{toc}{
    name={translocon on the outer chloroplast membrane},
    description={Proteinkomplex in der äußeren Chloroplastenmembran, erkennt Transitpeptide kernkodierter Vorläuferproteine und leitet sie an TIC weiter.},
    type=defi
}

\newglossaryentry{tim}{
    name={translocator of the
    inner (mitochondrial)
    membrane},
    description={Proteinkomplex in der inneren Mitochondrienmembran, importiert kernkodierte Vorläuferproteine in die mitochondriale Matrix oder inseriert sie in die Innenmembran.},
    type=defi
}

\newglossaryentry{tom}{
    name={translocator of the
    outer (mitochondrial)
    membrane},
    description={Proteinkomplex in der äußeren Mitochondrienmembran, erste Erkennungs- und Translokationsstelle für kernkodierte mitochondriale Vorläuferproteine.},
    type=defi
}

\newglossaryentry{transkription}{
    name={Transkription},
    description={Die Synthese von RNA anhand einer DNA-Matrize},
    type=defi
}

\newglossaryentry{trailer}{
    name={Trailer},
    description={3'-nicht-translatierter Bereich (3'-UTR) einer mRNA, der nach dem Stoppcodon liegt und regulatorische Sequenzen wie das Polyadenylierungssignal enthalten kann.},
    type=defi
}

\newglossaryentry{topoisomerase}{
    name={Topoisomerase},
    description={Ein Enzym, das die Anzahl der Kreuzungspunkte der beiden Stränge
    in einem geschlossenen DNA-Molekül verändert. Dazu
    spaltet es die DNA, führt die DNA durch den Bruch und verbindet
    die DNA wieder miteinander. (Lewin)},
    type=defi
}

\newglossaryentry{tf}{
    name={Terminationsfaktor},
    description={Protein oder Proteinkomplex, der die Beendigung der Transkription vermittelt, indem er die RNA-Polymerase vom DNA-Template und das RNA-Transkript freisetzt.},
    type=defi
}

\newglossaryentry{telomerase}{
    name={Telomerase},
    description={Das Ribonukleoprotein-Enzym, das am Telomer sich wiederholende Einheiten aus
    einem Strang bildet, indem es einzelne Basen an das
    3'-Ende der DNA anfügt, gemäß der Vorgabe einer RNA-Sequenz in der RNA-Komponente
    des Enzyms. (Lewin)},
    type=defi
}

\newglossaryentry{telomer}{
    name={Telomer},
    description={Die Enden der Chromosomen},
    type=defi
}

\newglossaryentry{tgs}{
    name={transcriptional gene silencing},
    description={epigenetische Mechanismen der Transkriptionsregulation},
    type=defi
}

\newglossaryentry{translocon}{
    name={Translocon},
    description={Proteinkanal in biologischen Membranen, dessen Funktion darin besteht, Polypeptide durch die Membran zu
    transportieren oder sie in die Lipid-Doppelschicht einzubetten},
    type=defi
}

\newglossaryentry{trna}{
    name={tRNA},
    first={transfer RNA (tRNA)},
    description={NA-Moleküle, die wäh-
    rend der Translation spezifische Aminosäuren zu den
    Ribosomen transportieren},
    type=defi
}

\newglossaryentry{spleißosom}{
    name={Spleißosom},
    plural={Spleißosomen},
    description={Organell (Großer Ribonukleoprotein-Komplex aus snRNAs und Proteinen), das das Spleißen
    der prä-mRNA katalysiert, indem Introns entfernt und Exons miteinander verknüpft werden.},
    type=defi
}


\newglossaryentry{superfamilie}{
    name={Superfamilie},
    description={Eine Gruppe von Genen, die alle durch eine vermutete Abstammung von einem gemeinsamen
    Vorfahren miteinander verbunden sind, mittlerweile jedoch erhebliche Unterschiede aufweisen. (Lewin)},
    type=defi
}


\newglossaryentry{syntenie}{
    name={Syntenie},
    description={Physikalische Ko-Lokalisierung von genetischen Loci (Reihenfolge
    von Loci) auf äquivalenten Chromosomen verschiedener Spezies},
    type=defi
}


\newglossaryentry{ssr}{
    first={simple sequence repeat (SSR)},
    name={SSR},
    description={= Mikrosatellit, Minisatellit, STR},
    type=defi
}

\newglossaryentry{spleißen}{
    name={Spleißen},
    description={Entfernen der Introns},
    type=defi
}

\newglossaryentry{snrnp}{
    name={snRNP},
    first={kleiner nukleärer Ribonukleoprotein-Komplex (snRNP)},
    description={Komplex aus kleiner nukleärer RNA (snRNA) und Proteinen. Hauptbestandteil des Spleißosoms, katalysiert das Spleißen von prä-mRNA.},
    type=defi
}

\newglossaryentry{snorna}{
    name={snoRNA (small nucleolar RNA)},
    plural={snoRNAs},
    description={Kleine, im Nukleolus lokalisierte RNAs, die an der Modifikation und Prozessierung der rRNA (endonukleolytische Spaltung, Nukleotid-Modifikation) beteiligt sind.},
    type=defi
}


\newglossaryentry{silencing}{
    name={Silencing},
    description={Sequenzspezifische Repression der Genexpression auf
    transkriptioneller (TGS) oder post-transkriptioneller Ebene (PTGS).
    Mechanismen umfassen Chromatinkondensation, DNA-Methylierung sowie
    gezielte Degradierung oder Translationshemmung spezifischer mRNAs,
    z.\,B.\ durch RNA-Interferenz.},
    type=defi
}


\newglossaryentry{sine}{
    name={SINE},
    first={SINE (short-interspersed nuclear element)},
    description={Eine bedeutende Klasse kurzer (weniger als 500 bp) nicht-autonomer
    Retrotransposons, die etwa 13 \% des menschlichen
    Genoms einnehmen. (Lewin)},
    type=defi
}

\newglossaryentry{silencer}{
    name={Silencer},
    plural={Silencer},
    description={Distales regulatorisches DNA-Element, das die Transkription eines Gens durch Bindung repressiver Transkriptionsfaktoren hemmt.},
    type=defi
}

\newglossaryentry{srp}{
    name={Signal Recognition Particle},
    first={Signal Recognition Particle (SRP)},
    description={Ribonukleoprotein, das am cotranslationalen Transport von Proteinen in das endoplasmatische Retikulum (ER)
    von Eukaryoten und die Plasmamembran von Prokaryoten beteiligt ist},
    type=defi
}

\newglossaryentry{signalpeptidase}{
    name={Signal-Peptidase},
    description={Proteinspaltendes Enzym (Protease), die an der Membran des Endoplasmatischen Retikulums (ER) innerhalb
    der Zelle lokalisiert sind. Erkennt spezielle Aminosäuresequenzen (Signalsequenzen) innerhalb eines Proteins. Sie
    spaltet im Lumen des ER das Signalpeptid von der restlichen Polypeptidkette ab},
    type=defi
}

\newglossaryentry{snrna}{
    name={snRNA},
    description={Small nuclear RNA; Bestandteil des Spleißosoms und an der Prozessierung von prä-mRNA durch Spleißen beteiligt.},
    type=defi
}

\newglossaryentry{sncrna}{
    name={sncRNA},
    description={Small non-coding RNA; Sammelbegriff für kurze nicht-kodierende RNAs, die regulatorische Funktionen in der Genexpression übernehmen.},
    type=defi
}

\newglossaryentry{sirna}{
    name={siRNA},
    description={Small interfering RNA; kurze regulatorische RNA (ca. 21–24 Nukleotide), die durch RNA-Interferenz spezifische mRNAs abbaut oder deren Translation hemmt.},
    type=defi
}

\newglossaryentry{sigma_faktor}{
    name={Sigma-Faktor},
    description={Unterheit der bakteriellen RNA-Polymerase, die die spezifische Erkennung von Promotorsequenzen ermöglicht und die Initiation der Transkription steuert.},
    type=defi
}
\newglossaryentry{semiRep}{
    name={semikonservative Replikation},
    description={Tatsächlicher Replikationsmodus der DNA, bei dem jede Tochter-Doppelhelix aus einem alten Matrizenstrang und einem neu synthetisierten Strang besteht (Meselson-Stahl-Experiment).},
    type=defi
}


\newglossaryentry{rnapolv}{
    name={RNA-Polymerase V},
    description={ranskribiert ``intergenic long ncRNA''},
    type=defi
}

\newglossaryentry{rnapolii}{
    name={RNA-Polymerase II},
    description={Im Nukleoplasma lokalisierte RNA-Polymerase, die die Synthese der 5S‑rRNAs, tRNAs, snRNAs und weiterer
    kleiner RNAs katalysiert},
    type=defi
}
\newglossaryentry{rnapoliii}{
    name={RNA-Polymerase III},
    description={Im Nukleoplasma lokalisierte RNA-Polymerase, die für die Synthese der mRNAs bzw. ihrer Vorläufer, die
    hnRNAs, sowie miRNAs zuständig ist},
    type=defi
}


\newglossaryentry{rnapoli}{
    name={RNA-Polymerase I},
    description={Im Nukleolus lokalisierte RNA-Polymerase, die für die Synthese der rRNAs (außer der 5S‑rRNA) verantwortlich ist},
    type=defi
}
\newglossaryentry{rnapol}{
    name={RNA-Polymerase},
    description={Enzym, das RNA anhand einer DNA- (oder RNA-) Matrize synthetisiert. In Eukaryoten existieren mehrere Typen (Pol I–III), die unterschiedliche RNA-Klassen transkribieren.},
    type=defi
}

\newglossaryentry{rnapoliv}{
    name={RNA-Polymerase IV},
    description={Pflanzenspezifische RNA-Polymerase, die an der Bildung von siRNAs und der RNA-gerichteten DNA-Methylierung beteiligt ist.},
    type=defi
}


\newglossaryentry{rna}{
    name={RNA},
    description={Ribonukleinsäure; ein meist einzelsträngiges Nukleinsäuremolekül, das aus Ribonukleotiden besteht und zentrale Funktionen in der Genexpression, Katalyse und Regulation übernimmt.},
    type=defi
}


\newglossaryentry{rdrp}{
    name={RNA-dependent RNA polymerase (RDRP)},
    description={RNA-abhängige RNA-Polymerase, die RNA als Matrize verwendet, um neue RNA-Stränge zu synthetisieren. Kommt z. B. bei RNA-Viren und in bestimmten RNA-Interferenz-Systemen (z. B. siRNA-Verstärkung bei Pflanzen) vor.},
    type=defi
}

\newglossaryentry{rnai}{
    name={RNA-Interferenz (RNAi)},
    description={Post-transkriptioneller Gen-Silencing-Mechanismus, bei dem
    kurze doppelsträngige RNAs (siRNA, miRNA) nach Prozessierung durch
    Dicer in den RISC-Komplex (RNA-induced silencing complex    )
    geladen werden. RISC erkennt komplementäre mRNAs und bewirkt deren
    Degradierung (vollständige Komplementarität) oder Translationsrepression
    (partielle Komplementarität).},
    type=defi
}

\newglossaryentry{risc}{
    name={RISC-Komplex (RNA-induced silencing complex)},
    description={Effektorkomplex der
    RNA-Interferenz. Kernkomponente ist ein Argonaute-Protein (AGO),
    das einen einzelsträngigen Guide-RNA-Strang (siRNA oder miRNA) bindet
    und komplementäre Ziel-mRNA erkennt. Bei vollständiger Komplementarität
    erfolgt endonukleolytische Spaltung (Slicer-Aktivität von AGO2);
    bei partieller Komplementarität wird die Translation reprimiert
    oder die mRNA destabilisiert.},
    type=defi
}

\newglossaryentry{restriktionskarte}{
    name={Restriktionskarte},
    description={eine lineare Abfolge von Stellen auf der DNA, die durch definierte Abstände voneinander getrennt sind bzw. die
    Darstellung einer linearen Abfolge der Stellen, an denen bestimmte Restriktionsenzyme ihre Zielstellen finden.
    Wird ein DNA-Molekül mit einem geeigneten Restriktionsenzym geschnitten, wird es in einzelne, negativ geladene Fragmente gespalten.
    Diese Fragmente lassen sich anhand ihrer Größe durch Gelelektrophorese trennen. Durch die Analyse der Restriktionsfragmente der DNA ist es möglich,
    eine Karte des ursprünglichen Moleküls zu erstellen.
    },
    type=defi
}
% Lewin

\newglossaryentry{regulatorische_rna}{
    name={regulatorische RNA},
    plural={regulatorische RNAs},
    description={Sammelbegriff für nicht-kodierende RNAs, die die Genexpression auf Transkriptions-,
    posttranskriptioneller oder translationaler Ebene regulieren, z. B. miRNA, siRNA oder lncRNA.},
    type=defi
}

\newglossaryentry{redundantes_gen}{
    name={redundantes Gen},
    description={Der Phänotyp einer Mutante mit einem Defekt in
    solchen Genen ist nicht unterscheidbar vom Wildtyp. (Vorlesung 3)},
    type=defi
}

\newglossaryentry{replikon}{
    name={Replikon},
    description={Replikationseinheit der
    Erbsubstanz, die einen Startpunkt der Replikation enthält und zur autonomen Replikation befähigt ist. (Allgemeine Genetik)},
    type=defi
}


\newglossaryentry{retrotransposon}{
    name={Retrotransposon},
    description={Ein Transposon, das über eine RNA-Form mobilisiert; das DNA-Element
    wird in RNA transkribiert und anschließend in
    DNA revers transkribiert, die an einer neuen Stelle im Genom inseriert wird. Es besitzt keine
    infektiöse (virale) Form.},
    type=defi
}

\newglossaryentry{rdna}{
    name={rDNA},
    first={ribosomale DNA (rDNA)},
    description={Enthält Gene für die für rRNA (rRNA-Gene)},
    type=defi
}
\newglossaryentry{rbs}{
    name={RBS},
    first={Shine-Dalgarno-Sequenzen (RBS)},
    description={Purinreiche Sequenz in prokaryotischen mRNAs (Konsensus: 5'-AGGAGG-3'), die ca. 5--10\,nt
    upstream des Startcodons liegt und durch komplementäre Basenpaarung mit der 16S-rRNA die Ribosomenbindung vermittelt.},
    type=defi
}

\newglossaryentry{pseudogen}{
    name={Pseudogen},
    description={Gene, die ihre Fähigkeit verloren haben,
    funktionsfähige Versionen der Polypeptide zu produzieren, für die sie ursprünglich
    kodierten. Die Gründe für ihre Inaktivität sind vielfältig: Die Defizite
    können bei der Transkription, der Translation oder bei beiden liegen. Sie können nicht funktionsfähig sein oder eine veränderte
    Funktion aufweisen, und die RNA-Produkte könnten regulatorische Funktionen erfüllen. (Lewin)},
    type=defi
}

\newglossaryentry{processed_pseudogen}{
    name={prozessierte Pseudogene},
    description={Pseudogene, die durch reverse
    Transkription und Integration von mRNA-Transkripten entstehen.},
    type=defi
}


\newglossaryentry{pic}{
    name={Prä-Initiationskomplex},
    first={Prä-Initiationskomplex (Pre-Initiation Complex, PIC)},
    plural={Prä-Initiationskomplexe (Pre-Initiation Complexes, PICs)},
    description={Multiproteinkomplex aus RNA-Polymerase II und allgemeinen Transkriptionsfaktoren (TFIID, TFIIA, TFIIB, TFIIF, TFIIE, TFIIH), der sich am Core-Promotor bildet und die Initiation der Transkription ermöglicht.},
    type=defi
}


\newglossaryentry{promotor}{
    name={Promotor},
    plural={Promotoren},
    description={DNA-Region vor einem Gen, die die Bindung der RNA-Polymerase und der Transkriptionsfaktoren ermöglicht und den Startpunkt der Transkription festlegt.},
    type=defi
}

\newglossaryentry{prämrna}{
    name={prä-mRNA},
    description={Unreifes, primäres RNA-Pol-II-Transkript im Zellkern
    vor der Prozessierung. Enthält Introns und Exons sowie nicht-translatierte
    Bereiche. Wird durch 5'-Capping, Spleißen (Entfernung der Introns)
    und 3'-Polyadenylierung zur reifen mRNA prozessiert,
    die anschließend ins Zytoplasma exportiert wird.},
    type=defi
}

\newglossaryentry{premirna}{
    name={pre-miRNA},
    description={Ca.\ 60--70\,nt lange Haarnadel-RNA, die durch Drosha-
    Prozessierung aus der pri-miRNA entsteht. Wird durch Exportin-5
    in das Zytoplasma transportiert und dort von Dicer zur reifen
    miRNA prozessiert.},
    type=defi
}


\newglossaryentry{primimrna}{
    name={pri-miRNA},
    description={Primäres RNA-Pol-II-Transkript, das eine oder mehrere
    Haarnadelstrukturen (Hairpins) enthält. Wird im Zellkern durch den
    Microprocessor-Komplex (Drosha\,+\,DGCR8) zur ca.\ 60--70\,nt
    langen pre-miRNA prozessiert.},
    type=defi
}

\newglossaryentry{protein_kinase}{
    name={Protein-Kinase},
    plural={Protein-Kinasen},
    description={Enzym, das Proteine durch Übertragung einer Phosphatgruppe (meist von ATP auf Serin-, Threonin- oder Tyrosinreste) phosphoryliert und dadurch deren Aktivität, Interaktionen oder Lokalisation reguliert.},
    type=defi
}

\newglossaryentry{primaertranskript}{
    name={Primärtranskript},
    plural={Primärtranskripte},
    description={Das unmittelbare Produkt der Transkription eines Gens. Bei eukaryotischen proteincodierenden Genen entspricht es der prä-mRNA (hnRNA) und wird durch Prozessierungsschritte wie 5'-Capping, Spleißen und Polyadenylierung zur reifen mRNA verarbeitet.},
    type=defi
}

\newglossaryentry{ptgs}{
    name={post-transkriptionelles Gene Silencing (PTGS)},
    description={Regulation der Genexpression auf mRNA-Ebene nach der
    Transkription; umfasst zwei Hauptmechanismen: (1)~gezielte Degradierung
    spezifischer mRNAs und (2)~gezielte Inhibierung ihrer Translation.
    Wird u.\,a.\ durch RNA-Interferenz (siRNA, miRNA) vermittelt.
    Bedeutend für endogene Genregulation sowie als Werkzeug der reversen
    Genetik zur Analyse von Genfunktionen (Knockdown).},
    type=defi
}


\newglossaryentry{pas}{
    name={Polyadenylierungssignal},
    description={Sequenz im 3'-UTR eukaryotischer prä-mRNA (meist AAUAAA), die die Spaltung der RNA und anschließende Polyadenylierung durch die Poly(A)-Polymerase auslöst.},
    type=defi
}

\newglossaryentry{polya_polymerase}{
    name={Poly(A)-Polymerase},
    plural={Poly(A)-Polymerasen},
    description={Enzym, das nach der Transkriptionsspaltung am 3'-Ende eukaryotischer prä-mRNA eine Poly(A)-Sequenz synthetisiert, indem es Adenosinreste ohne DNA-Matrize anfügt.},
    type=defi
}

\newglossaryentry{polya}{
    name={Poly-A-Schwanz},
    description={Posttranskriptionelle Modifikation am 3'-Ende eukaryotischer mRNA, bestehend aus einer Kette von Adenin-Nukleotiden, die Stabilität, Export aus dem Zellkern und Translationseffizienz erhöht.},
    type=defi
}

\newglossaryentry{polycistronisch}{
    name={polycistronisch},
    description={Kodierend für mehrere Polypedtide},
    type=defi
}

\newglossaryentry{plastom}{
    name={Plastom},
    description={Die Erbinformationen aus den Chlorpolasten; das Chlorpolasten-Genom},
    type=defi
}


\newglossaryentry{paralog}{
    name={Paralogie},
    first={paraloge Gene},
    description={homologe Gene in einer Art, die durch Duplikation
    entstanden sind},
    type=defi
}

\newglossaryentry{patVer}{
    name={paternale Vererbung},
    description={Vererbung genetischer Information über den väterlichen Elternteil},
    type=defi
}


\newglossaryentry{ortholog}{
    name={Orthologie},
    first={orthologe Gene},
    description={Gene in unterschiedlichen Arten, die von einer gemeinsamen
    Sequenz abstammen. Sie können die gleiche Funktion haben, müssen es aber nicht.},
    type=defi
}
\newglossaryentry{orf}{
    name={Open Reading Frame (ORF)},
    plural={Open Reading Frames (ORFs)},
    description={Leserahmen einer mRNA, der zwischen Startcodon (AUG) und Stopcodon liegt und die Information für die Proteinsequenz enthält.},
    type=defi
}

\newglossaryentry{operon}{
    name={Operon},
    description={Genetische Funktionseinheit bei Bakterien, die aus einer Gruppe von Genen besteht, die
    über einen gemeinsamen Promotor reguliert werden und
    eine gemeinsame mRNA (Messenger) produzieren, die
    für mehrere Proteine kodiert. (Allgemeine Genetik)},
    type=defi
}


\newglossaryentry{nukleosom}{
    name={Nukleosom},
    description={Grundlegende Verpackungseinheit des eukaryotischen Chromatins: ca. 147\,bp DNA sind 1{,}65-mal linkshändig um einen Histonoktamer (je zwei Kopien von H2A, H2B, H3 und H4) gewickelt},
    type=defi
}
\newglossaryentry{nukleolus}{
    name={Nukleolus},
    plural={Nukleoli},
    description={Struktur im Zellkern, in der die rRNA-Synthese und der Zusammenbau der ribosomalen Untereinheiten stattfindet; bildet sich an den Nukleolus-Organisator-Regionen (NORs).},
    type=defi
}

\newglossaryentry{nukleotid}{
    name={Nukleotid},
    description={Baustein der Nucleinsäuren, der aus einem Nucleosid und einer
    daran gebundenen Phosphatgruppe besteht},
    type=defi
}

\newglossaryentry{nukleosid}{
    name={Nukleosid},
    description={Bausteine der Nukleinsäuren, die aus je einer Purin- oder Pyrimidinbase und einem Ribose- oder Desoxyribosezucker bestehen},
    type=defi
}

\newglossaryentry{nor}{
    name={Nukleolusorganisatorregion},
    description={Chromosomale DNA-Region, die rRNA-Gene enthält und als Organisationszentrum für den Nukleolus dient, wo die rRNA-Synthese und Ribosomenassemblierung stattfinden.},
    type=defi
}

\newglossaryentry{nukleoporin}{
    name={Nukleoporin},
    description={Eine Proteinfamilie, aus der der NPC aufgabaut ist},
    type=defi
}

\newglossaryentry{nukleom}{
    name={nucleares Genom},
    description={Die Erbinformationen aus dem Zellkern.},
    type=defi
}

\newglossaryentry{nls}{
    name={Nuclear Localization Signal},
    description={Kurze, meist basische Aminosäuresequenz in einem Protein, die von Importinen erkannt wird und den Kernimport über NPCs vermittelt.},
    type=defi
}

\newglossaryentry{nes}{
    name={Nuclear Export Signal},
    description={Kurze, meist leucinreiche Aminosäuresequenz, die von Exportinen (z. B. CRM1) erkannt wird und den Export eines Proteins aus dem Zellkern vermittelt.},
    type=defi
}

\newglossaryentry{ndna}{
    name={nDNA},
    description={Nukleäre DNA; die im Zellkern lokalisierte Hauptform des eukaryotischen Genoms.},
    type=defi
}

\newglossaryentry{sense}{
    name={nicht-kodogener Strang},
    plural={nicht-kodogene Stränge},
    description={DNA-Strang, der die gleiche Basensequenz wie die transkribierte mRNA (bis auf Uracil statt Thymin) besitzt und nicht als Vorlage für die RNA-Synthese dient; auch Plus-Strang oder codierender Strang genannt.},
    type=defi
}

\newglossaryentry{nonprocessed_pseudogen}{
    name={nicht-prozessierte Pseudogene},
    description={Pseudogene, die durch unvollständige
    Duplikation oder Mutation der zweiten Kopie funktioneller Gene entstehen.},
    type=defi
}


\newglossaryentry{ncrna}{
    name={ncRNA},
    description={Non-coding RNA; Sammelbegriff für alle RNAs, die nicht in Proteine übersetzt werden, aber strukturelle, katalytische oder regulatorische Funktionen besitzen.},
    type=defi
}

\newglossaryentry{nterminus}{
    name={N-terminus},
    description={Aminogruppe am Anfang eines Proteins},
    type=defi
}


\newglossaryentry{mutante}{
    name={Mutante},
    description={Ein durch eine Mutation neu entstandenes Merkmal},
    type=defi
}


\newglossaryentry{mutagenese}{
    name={Mutagenese},
    description={Entstehung oder gezielte Herbeiführung von Mutationen im Erbgut durch endogene (Replikationsfehler) oder exogene Einwirkung (Mutagene).},
    type=defi
}

\newglossaryentry{mtdna}{
    name={mtDNA},
    description={Mitochondriale DNA; ringförmiges Genom der Mitochondrien, das für einen Teil der mitochondrialen Proteine sowie rRNAs und tRNAs kodiert.},
    type=defi
}

\newglossaryentry{mrna}{
    name={mRNA},
    description={Messenger-RNA; RNA, die die genetische Information von der DNA zu den Ribosomen überträgt und als Vorlage für die Proteinsynthese dient.},
    type=defi
}
\newglossaryentry{mrnp}{
    name={mRNP},
    first={Messenger-Ribonukleoprotein-Komplex (mRNP)},
    description={Komplex aus reifer mRNA und gebundenen Proteinen. In dieser Form wird die mRNA aus dem Kern exportiert und im Zytoplasma transportiert und translatiert.   },
    type=defi
}

\newglossaryentry{monocistronisch}{
    name={monocistronisch},
    description={nur für ein Polypeptid kodierend},
    type=defi
}

\newglossaryentry{minisatellit}{
    name={Minisatellit},
    description={Abschnitte der DNA im Genom bezeichnet, die aus tandemartigen Wiederholungen einer kurzen (ca. 12–100 Nukleotiden langer) DNA-Sequenz bestehen.; 10–60bp},
    type=defi
}
%  \cite{dewiki:Minisatellit, Knust2025_Modellorganismen}


\newglossaryentry{mirna}{
    name={miRNA (microRNA)},
    description={MicroRNA; kurze regulatorische RNA (ca. 20–22 Nukleotide), die durch Basenpaarung mit Ziel-mRNAs deren Translation hemmt oder deren Abbau fördert.},
    type=defi
}

\newglossaryentry{mirnaduplex}{
    name={miRNA-miRNA* Duplex},
    description={Kurzlebiges, ca.\ 21--23\,bp langes RNA-Duplex-Intermediat,
    das nach Dicer-Prozessierung der pre-miRNA entsteht. Der funktionelle
    Guide-Strang (miRNA) wird in RISC geladen; der komplementäre
    Passenger-Strang (miRNA*, auch miRNA-3p) wird meist degradiert.
    Die Strangauswahl erfolgt anhand der thermodynamischen Stabilität
    der Duplex-Enden: der Strang mit dem weniger stabilen 5'-Ende
    wird bevorzugt als Guide-Strang selektiert.},
    type=defi
}

\newglossaryentry{mikrosatellit}{
    name={Mikrosatellit},
    description={kurze wiederholende Abschnitte; 1–6bp; multi‐allelische polymorphe Marker; },
    type=defi
}

\newglossaryentry{matVer}{
    name={maternale Vererbung},
    description={Vererbung genetischer Information über den mütterlichen Elternteil},
    type=defi
}


\newglossaryentry{antisense}{
    name={Matrizenstrang},
    plural={Matrizenstränge},
    description={DNA-Strang, der als Vorlage (Template-Strang) für die Transkription dient und komplementär zur entstehenden mRNA ist; auch Antisense-, kodogener oder Template-Strang genannt.},
    type=defi
}


\newglossaryentry{leader}{
    name={Leader},
    description={5'-nicht-translatierter Bereich (5'-UTR) einer mRNA, der vor dem Startcodon liegt und an der Regulation von Translation und Stabilität beteiligt ist.},
    type=defi
}

\newglossaryentry{linkerdna}{
    name={Linker-DNA},
    description={DNA-Abschnitt (10--80\,bp) zwischen zwei aufeinanderfolgenden
    Nukleosomen; frei zugänglich für Nukleasen (z.\,B.\ MNase) und
    Transkriptionsfaktoren. Wird durch Histon H1 stabilisiert und
    beim vollständigen MNase-Verdau als erstes abgebaut.},
    type=defi
}


\newglossaryentry{line}{
    name={LINE},
    first={LINE (long-interspersed nuclear element)},
    description={Eine wichtige Klasse von Retrotransposons, die etwa
    21\% des menschlichen Genoms ausmachen. (Lewin)},
    type=defi
}


\newglossaryentry{kopplung}{
    name={Kopplung},
    description={Assoziation von Genen auf dem
    gleichen Chromosom, die zur gemeinsamen Vererbung der
    entsprechenden Merkmale führt},
    type=defi
}

\newglossaryentry{consensus_promoter}{
    name={Konsensus-Promotor},
    description={Idealisierte DNA-Sequenz, die aus der Auswertung vieler Promotoren abgeleitet wird und die am häufigsten vorkommenden Basen an jeder Position darstellt.},
    type=defi
}

\newglossaryentry{npc}{
    name={Kernporenkomplex},
    first={Kernporenkomplex (NPC)},
    description={Großer Proteinkomplex in der Kernhülle, der den selektiven Transport von Molekülen zwischen Zellkern und Zytoplasma ermöglicht.},
    type=defi
}
\newglossaryentry{karyotyp}{
    name={Karyotyp},
    description={Vollständiger Satz von Metaphase-
    chromosomen},
    type=defi
}
\newglossaryentry{inr}{
    name={Initiator-Element (INR)},
    description={Kernelement des RNA-Polymerase-II-Core-Promotors, das den Transkriptionsstart (+1) überlappt und die Initiation der Transkription unterstützt},
    type=defi
}

\newglossaryentry{internerpromotor}{
    name={interner Promotor},
    description={Promotor, dessen regulatorische Sequenzen innerhalb des transkribierten Gens liegen},
    type=defi
}

\newglossaryentry{invitro}{
    name={in vitro},
    description={Experimentelle Bedingungen außerhalb eines lebenden Organismus, z. B. in einem Reagenzglas oder unter definierten Laborbedingungen.},
    type=defi
}

\newglossaryentry{invivo}{
    name={in vivo},
    description={Vorgänge innerhalb eines lebenden Organismus unter physiologischen Bedingungen.},
    type=defi
}

\newglossaryentry{intron}{
    name={Intron},
    description={Abschnitt eines gespaltenen
    Gens, der beider Reifung der RNA  (Verspleißen)
    herausgeschnitten und demzufolge nicht in Protein übersetzt wird},
    type=defi
}

\newglossaryentry{hnrnp}{
    name={hRNP},
    first={heterogener Ribonukleoprotein-Komplex (mRNP)},
    description={Komplex aus heterogener nukleärer RNA (prä-mRNA) und Proteinen. Bindet naszierende Transkripte und ist an Prozessierung, Stabilität und Export der RNA beteiligt.},
    type=defi
}

\newglossaryentry{homolog}{
    name={Homologie},
    description={Ähnlichkeit biologischer Strukturen oder Sequenzen aufgrund gemeinsamer evolutionärer Abstammung},
    type=defi
}

\newglossaryentry{hnrna}{
    name={hnRNA},
    first={heterogene nukleäre RNA (hnRNA)},
    description={RNA-Klasse heterogener Größe, die im Kern lokalisiert ist und zu der vor allem die Vorläufer der mRNAs gehören, die im Kern in unreifer Form als Primärtranskripte vorliegen},
    type=defi
}


\newglossaryentry{histon}{
    name={Histon},
    description={Kleine, basische Kernproteine (reich an Lysin und Arginin),
    die den Histonoktamer des Nukleosoms bilden.},
    type=defi
}

\newglossaryentry{heterosom}{
    name={Heterosom},
    description={Geschlechtschromosom (Allosom), das sich strukturell oder im Replikationsverhalten vom Autosomensatz unterscheidet},
    type=defi
}

\newglossaryentry{helikase}{
    name={Helikase},
    plural={Helikasen},
    description={Enzym, das Nukleinsäure-Doppelstränge (DNA oder RNA) unter ATP-Verbrauch entwinden und trennen kann, z. B. während Replikation oder Transkription.},
    type=defi
}

\newglossaryentry{heterochromatin}{
    name={Heterochromatin},
    description={hochverpackte, nicht zugängliche DNA in
    Metaphasechromosomen oder selektiv kondensierten
    Genomabschnitten;
    ist Transkriptions-inaktiv und wird spät repliziert},
    type=defi
}

\newglossaryentry{hairpin_rna}{
    name={Hairpin-Struktur (RNA)},
    plural={Hairpin-Strukturen (RNA)},
    description={Sekundärstruktur der RNA, bei der sich eine einzelsträngige RNA durch intramolekulare Basenpaarung zu einer Schleifenstruktur (Stem-Loop) faltet; häufig an der Regulation und Prozessierung beteiligt.},
    type=defi
}

\newglossaryentry{gen}{
    name={gen},
    description={unit of hereditary information},
    type=defi
}
\newglossaryentry{genom}{
    name={Genom},
    description={Alle Erbinformationen (DNA-Sequenzen, Gene), inkl. der aus bspw. den Mitochondrien.},
    type=defi
}

\newglossaryentry{genet-marker}{
    name={genetischer Marker},
    description={Ein Gen oder eine DNA-Sequenz mit bekannter Position auf einem Chromosom, das bzw. die zur Identifizierung von Individuen oder Arten verwendet werden kann / Eine erkennbare DNA-Variation an einer bestimmten Stelle im Genom},
    type=defi
}

\newglossaryentry{gef}{
    name={guanine nucleotude exchange factor},
    description={Proteine oder Proteindomänen, die monomere GTPasen aktivieren, indem sie die Freisetzung von
    Guanosindiphosphat (GDP) anregen, um die Bindung von Guanosintriphosphat (GTP) zu ermöglichen},
    type=defi
}

\newglossaryentry{guiderna}{
    name={Guide-RNA},
    description={Einzelsträngige RNA, die einen Effektorkomplex durch
    Basenpaarung zur komplementären Zielsequenz dirigiert. Im RNAi-Kontext:
    der in RISC geladene miRNA- oder siRNA-Strang; im CRISPR-Kontext:
    die sgRNA, die Cas9 zur Ziel-DNA leitet.},
    type=defi
}

\newglossaryentry{genfamilie}{
    name={Genfamilie},
    description={Eine Gruppe von Genen, die
    infolge von Genduplikationen
    verwandte oder identische Produkte kodieren.
    Nach dem ersten Duplikationsvorgang sind die beiden Kopien
    identisch, divergieren jedoch anschließend, da sich in ihnen unterschiedliche Mutationen
    ansammeln. Weitere Duplikationen und Divergenzen erweitern die Familie. (Lewin)},
    type=defi
}


\newglossaryentry{gap}{
    name={GTPase-activating protein},
    description={Eine Familie von regulatorischen Proteinen, deren Mitglieder an aktivierte G-Proteine binden und deren
    GTPase-Aktivität stimulieren können, wodurch der Signalweg beendet wird},
    type=defi
}

\newglossaryentry{gprotein}{
    name={G-Protein},
    description={ Eine Familie von Proteinen, die innerhalb von Zellen als molekulare Schalter fungieren und an der
    Weiterleitung von Signalen von verschiedenen Reizen außerhalb der Zelle in ihr Inneres beteiligt sind.
    Ihre Aktivität wird durch Faktoren reguliert, die ihre Fähigkeit steuern, an Guanosintriphosphat (GTP) zu binden und
    dieses zu Guanosindiphosphat (GDP) zu hydrolysieren. Wenn sie an GTP gebunden sind, befinden sie sich im
    ``Ein''-Zustand, und wenn sie an GDP gebunden sind, befinden sie sich im ``Aus''-Zustand. Sie gehören zur
    größeren Gruppe von Enzymen, die als GTPasen bezeichnet werden. },
    type=defi
}

\newglossaryentry{gcbox}{
    name={GC-Box},
    description={GC-reiches regulatorisches Promotorelement, das Bindestellen für spezifische Transkriptionsfaktoren enthält},
    type=defi
}


\newglossaryentry{gtpase}{
    name={GTPase},
    description={Enzym, das GTP zu GDP und anorganischem Phosphat hydrolysiert und dadurch als molekularer Schalter in Signaltransduktion und Transportprozessen fungiert.},
    type=defi
}


\newglossaryentry{ran}{
    name={GTPase Ran},
    first={Ras-related nuclear protein (Ran)},
    description={ Kleines G-Protein, das zwischen einer GTP- und einer GDP-gebundenen Form wechselt und dadurch als
    molekularer Schalter für den Kern-Zytoplasma-Transport dient. RanGTP ist im Zellkern angereichert, RanGDP im Zytoplasma,
    wodurch die Richtung des Transports durch NPCs festgelegt wird.},
    type=defi
}

\newglossaryentry{fgrepeat}{
    name={FG-Repeat},
    description={kleine hydrophobe Segmente, die lange Abschnitte hydrophiler Aminosäuren unterbrechen. Diese flexiblen
    Abschnitte bilden entfaltete oder ungeordnete Segmente ohne feste Struktur.[6] Sie bilden ein Kettengeflecht, durch das
    kleinere Moleküle diffundieren können, das jedoch große hydrophile Makromoleküle ausschließt. Diese großen Moleküle können
    eine Kernpore nur passieren, wenn sie von einem Signalmolekül begleitet werden, das vorübergehend mit dem FG-Wiederholungssegment
    eines Nukleoporins interagiert},
    type=defi
}



\newglossaryentry{exVer}{
    name={extranukleäre Vererbung},
    description={Vererbung genetischer Information außerhalb des Zellkerns, insbesondere über mitochondriale oder plastidäre DNA},
    type=defi
}


\newglossaryentry{exon}{
    name={Exon},
    description={Für RNA und Protein kodierender Ab-
    schnitt eines gespaltenen Gens. Bleibt beim
    Verspleißen erhalten},
    type=defi
}


\newglossaryentry{euchromatin}{
    name={Euchromatin},
    description={dekondensierte DNA (die aber immer noch an
    Nucleosomen gebunden ist); aktiv, hat transkribierte Regionen},
    type=defi
}


\newglossaryentry{essentielles_gen}{
    name={essentielles Gen},
    description={Der ``Phänotyp'' einer Mutante mit einem Defekt in
    solchen Genen ist ``nicht lebensfähig'' (lethal). (Vorlesung 3)},
    type=defi
}

\newglossaryentry{enhancer}{
    name={Enhancer},
    plural={Enhancer},
    description={Distales regulatorisches DNA-Element, das die Transkriptionsrate eines Gens erhöhen kann, indem es mit Promotorregionen über DNA-Schleifen interagiert.},
    type=defi
}


\newglossaryentry{endosymbionten}{
    name={Endosymbionten-Hypothese},
    description={Hypothese, dass Eukaryoten aus einer Endosymbiose prokaryotischer Vorläuferorganismen hervorgegangen sind. Demnach sind chemo- und phototrophe Bakterien von Archaeen aufgenommen worden, in denen sie sich zu Zellorganellen ihrer Wirtszellen entwickelt haben, darunter Mitochondrien und Plastiden.},
    type=defi
}

\newglossaryentry{endonuklease}{
    name={Endonuklease},
    description={Enzym, das Phosphodiesterbindungen innerhalb eines
    Nukleinsäurestrangs spaltet (im Gegensatz zur Exonuklease, die vom
    Ende her abbaut)},
    type=defi
}

\newglossaryentry{eef}{
    name={Elongationsfaktor (eEF)},
    plural={Elongationsfaktoren (eEFs)},
    description={Proteinfaktoren der eukaryotischen Translation, die den Ablauf der Elongation im Ribosom steuern, insbesondere die Anlieferung von Aminoacyl-tRNAs, die Peptidbindung und die Translokation des Ribosoms entlang der mRNA.},
    type=defi
}

\newglossaryentry{deadenylierung}{
    name={Deadenylierung},
    plural={Deadenylierungen},
    description={Enzymatischer Prozess der schrittweisen Verkürzung des Poly(A)-Schwanzes eukaryotischer mRNA, der häufig den Beginn des mRNA-Abbaus einleitet.},
    type=defi
}

\newglossaryentry{dnarnapol}{
    name={DNA-abhängige RNA-Polymerase},
    description={Enzym, das RNA anhand einer DNA-Matrize synthetisiert. In Eukaryoten existieren drei Haupttypen (RNA-Polymerase I, II und III), die jeweils unterschiedliche Klassen von Genen transkribieren.},
    type=defi
}

\newglossaryentry{cytogenetik}{
    name={Cytogenetik},
    description={todo},
    type=defi
}

\newglossaryentry{cyclin}{
    name={Cyclin},
    description={todo},
    type=defi
}

\newglossaryentry{cytokinese}{
    name={Cytokinese},
    description={Trennung der Zelle in zwei Tochterzellen nach der Mitose},
    type=defi
}

\newglossaryentry{cpsf}{
    name={CPSF (Cleavage and Polyadenylation Specificity Factor)},
    plural={CPSFs},
    description={Proteinkomplex, der die 3'-Endprozessierung eukaryotischer prä-mRNA erkennt und die Spaltung am Polyadenylierungssignal sowie die anschließende Polyadenylierung initiiert.},
    type=defi
}


\newglossaryentry{cpdna}{
    name={cpDNA},
    description={Chloroplasten-DNA; ringförmiges Genom der Chloroplasten, das für Gene der Photosynthese sowie eigene rRNAs und tRNAs kodiert.},
    type=defi
}

\newglossaryentry{corepromotor}{
    name={Core-Promotor},
    description={Minimaler DNA-Abschnitt eines Promotors, der für die Initiation der Transkription erforderlich ist},
    type=defi
}

\newglossaryentry{corenukleosom}{
    name={Core-Nukleosom},
    description={Minimale Nukleosomeneinheit nach Verdau der Linker-DNA:
    147\,bp DNA, die 1{,}65-mal linkshändig um den Core-Oktamer gewickelt
    sind. Enthält kein Histon H1; entspricht dem nucleosome core particle
    (NCP) in der Kristallstruktur.},
    type=defi
}

\newglossaryentry{coreoktamer}{
    name={Core-Oktamer},
    description={Proteinkern des Nukleosoms aus acht Core-Histonen:
    ein (H3$_2$--H4$_2$)-Tetramer als zentrale Grundstruktur, flankiert
    von zwei H2A--H2B-Dimeren. Hoch symmetrisch aufgebaut mit
    pseudo-dyadischer Symmetrieachse.},
    type=defi
}

\newglossaryentry{basaletrankskription}{
    name={basale Transkription},
    plural={basale Transkriptionen},
    description={Minimale Transkriptionsaktivität eines Gens, die allein durch den Core-Promotor, die allgemeinen Transkriptionsfaktoren und die RNA-Polymerase erreicht wird, ohne zusätzliche regulatorische Aktivatoren oder Repressoren.},
    type=defi
}

\newglossaryentry{autosom}{
    name={Autosom},
    description={Nicht-geschlechtsbestimmendes Chromosom; beim Menschen 22 homologe Paare, die in beiden Geschlechtern identisch vorliegen.},
    type=defi
}

\newglossaryentry{aminoacyl_trna}{
    name={Aminoacyl-tRNA},
    plural={Aminoacyl-tRNAs},
    description={Mit einer spezifischen Aminosäure beladene tRNA, die durch Aminoacyl-tRNA-Synthetasen erzeugt wird und die Aminosäure während der Translation zum Ribosom transportiert.},
    type=defi
}

\newglossaryentry{antisensetech}{
    name={Antisense-Technologie},
    description={konstitutive oder induzierte
    Expression von antisense RNA},
    symbol={Methode zur post-transkriptionellen Genexpressionsregulation
    durch einzelsträngige Antisense-Oligonukleotide (ASO), die komplementär
    zur Ziel-mRNA sind. Mechanismen: (1)~sterische Blockierung der
    Translation, (2)~RNase-H-vermittelte Degradierung des RNA:DNA-Hybrids.
    Im Gegensatz zu RNAi kein RISC-Komplex nötig; wichtiges Werkzeug
    für Knockdown-Analysen und therapeutische Anwendungen.},
    type=defi
}


\newglossaryentry{alu}{
    name={Alu-Repeats},
    description={Familie kurzer, etwa 300 bp langer repetitiver DNA-Sequenzen (SINEs), die im menschlichen Genom in sehr hoher Kopienzahl vorkommen.},
    type=defi
}


\newglossaryentry{allel}{
    name={Allel},
    description={variant form of a gene / Eine der verschiedenen Formen, die ein bestimmter Locus annehmen kann},
    type=defi
}


\newglossaryentry{dogma}{
    name={zentrales Dogma},
    description={Wenn (sequenzielle) Information einmal in ein Protein übersetzt wurde, kann sie dort nicht wieder
    herausgelangen},
    type=defi
}

\newglossaryentry{ab-initio}{
    name={``ab initio''-Genvorhersage},
    description={Computergestützte Vorhersage von Genen allein anhand der genomischen DNA-Sequenz und charakteristischer Sequenzmerkmale wie Start- und Stopcodons, Spleißsignalen, offenen Leserahmen oder Codon-Nutzung, ohne experimentelle Vergleichsdaten zu verwenden},
    type=defi
}

\newglossaryentry{zehnnmfaser}{
    name={10-nm-Faser (Nukleosomenfilament)},
    description={Erste Verdichtungsstufe des Chromatins; entspricht der
    Perlenketten-Struktur (``beads-on-a-string''): Nukleosomen sind
    durch Linker-DNA verbunden},
    type=defi
}

\newglossaryentry{chromosom}{
    name={Chromosom},
    description={Fädige Chromatinstruktur im Zellkern, Träger des genetischen
    Materials (DNA). Vor einer Zellteilung besteht ein C.
    aus 2 parallelen Chromatiden (2-Chromatid-Chromo-
    som), die am Centromer miteinander verbunden sind.},
    type=defi
}
\newglossaryentry{codon}{
    name={Codon},
    plural={Codons},
    description={Basentriplett in der mRNA, das für eine bestimmte Aminosäure oder ein Stop-Signal während der Translation codiert.},
    type=defi
}

\newglossaryentry{chromosomen_abberation}{
    name={Chromosomen-Abberation},
    description={Strukturelle oder numerische Veränderung von Chromosomen gegenüber dem Normalzustand, z. B. Deletion, Duplikation, Inversion, Translokation oder Aneuploidie.},
    type=defi
}


\newglossaryentry{chondrom}{
    name={Chondrom},
    description={Die Erbinformationen aus den Mitochondrien; das Mitochondrien-Genom},
    type=defi
}


\newglossaryentry{chiasma}{
    name={Chiasma},
    description={Die cytologisch
    sichtbare Folge von Crossover-
    Ereignissen.},
    type=defi
}

\newglossaryentry{chaperone}{
    name={Chaperone},
    description={ATPase, die die korrekte Faltung neu‐synthetisierter oder fehlgefalteter Proteine erleichtert},
    type=defi
}


\newglossaryentry{centromer}{
    name={Centromer},
    description={Spindelfaseransatzstelle (Spindel) der Chromosomen.},
    type=defi
}

\newglossaryentry{ctd}{
    name={Carboxyterminale Domäne (CTD)},
    plural={Carboxyterminale Domänen (CTDs)},
    description={Wiederholte Heptapeptid-Sequenz am C-Terminus der RNA-Polymerase II, die durch Phosphorylierung die
    Rekrutierung von Prozessierungsfaktoren und den Übergang von Initiation zu Elongation steuert.},
    type=defi
}


\newglossaryentry{cap}{
    name={Cap-Struktur},
    description={Modifizierte 5'-Endstruktur eukaryotischer mRNA (7-Methylguanosin-Kappe), die die mRNA vor Abbau schützt, den Export aus dem Zellkern ermöglicht und die Translation initiiert.},
    type=defi
}

\newglossaryentry{caatbox}{
    name={CAAT-Box},
    description={Konserviertes regulatorisches Promotorelement eukaryotischer Gene, das die Effizienz der Transkription beeinflusst},
    type=defi
}


\newglossaryentry{dreißignmfaser}{
    name={30-nm-Faser (30-nm Fibrille)},
    description={Zweite Verdichtungsstufe durch Aufwindung der 10-nm-Faser
    zu einer kompakteren Helix
    Histon H1 ist für die Stabilisierung essentiell.},
    type=defi
}

\newglossaryentry{fives_rrna}{
    name={5S-rRNA},
    plural={5S-rRNAs},
    description={Bestandteil der ribosomalen RNA der großen ribosomalen Untereinheit, der in Eukaryoten durch RNA-Polymerase III transkribiert wird.},
    type=defi
}


\newglossaryentry{bac}{
    name={BAC (Bacterial Artificial Chromosome)},
    description={Auf dem F-Plasmid basierender Klonierungsvektor, der sehr grosse DNA-Fragmente (ca. 100--300 kb) stabil in E. coli aufnehmen kann. Grundlage fuer physikalische Karten.},
    type=defi
}